\subsection{Integrals and integration systems}

\begin{definition}
  Let $X$ be a set and let $(X,L)$ be a Riesz system on $X$.
  We say that an $\mathbb{R}$-linear map
  \[
    I:L\to\mathbb{R}
  \]
  is an \textbf{integral} on the Riesz system $(X,L)$ if the following two
  conditions hold.
  \begin{itemize}
    \item For every $f\in L$, if $f\ge0$ then $I(f)\ge0$.
    \item For every sequence of functions $\{f_n\}\subset L$, if
      $f_n(x)\downarrow0$ for all $x\in X$ then $\lim_{n\to\infty}I(f_n)=0$.
  \end{itemize}
\end{definition}

\begin{definition}
  Let $(X,L)$ be a Riesz system on a set $X$ and let $I:L\to\mathbb{R}$
  be an integral on $(X,L)$.
  The triple
  \[
    (X,L,I)
  \]
  is called an \textbf{integration system} on the set $X$.
\end{definition}

\begin{example}[counting]\label{ex:counting}
  Let $X=\mathbb{N}$ and let $L$ be the set of all functions that vanish outside
  finitely many $n$.
  Clearly $L$ is a Riesz space on $X$.
  Putting
  \[
    I(f):=\sum_{n=1}^{\infty}f(n) ,
  \]
  the right-hand side is a finite sum and $I$ is an integral on $(X,L)$.
  Indeed, positivity is immediate, and if $f_n\downarrow0$ then the supports of
  the $f_n$ lie in the support of $f_1$, a finite set, so $I(f_n)\to0$.
\end{example}

\begin{example}[step functions and the Riemann integral]\label{ex:step-functions}
  Let $X=[0,1)$ and let $L$ be the set of all finite linear combinations of
  indicator functions of half-open intervals $[a,b)\subset[0,1)$, that is, the
  set of all step functions on $X$.
  Then $L$ is a Riesz space on $X$.
  Taking $I$ to be the Riemann integral
  \[
    I(f):=\int_0^1f(x)\,dx ,
  \]
  the triple $(X,L,I)$ is an integration system.

  Positivity is immediate. Continuity, that is, $I(f_n)\to0$ whenever
  $f_n(x)\downarrow0$ holds for every $x\in X$, follows from additivity of length
  over countable partitions of an interval.
  For a proof see \cite{BM}, \S2.
  The Daniell extension of this example is described in
  Example~\ref{ex:lebesgue}.
\end{example}

\begin{remark}\label{rem:why-half-open}
  There is a reason for taking the base space to be $[0,1)$ rather than $[0,1]$.
  If one takes $X=[0,1]$ and builds $L$ from indicator functions of half-open
  intervals, then every element of $L$ vanishes at $x=1$.
  Then $\sum_n\varphi_n(1)=0<+\infty$ for every series representation, so
  condition (3) of a series representation is imposed at $x=1$ always, and every
  element of $L^1$ satisfies $f(1)=0$.
  This differs from the set of all Lebesgue integrable functions on $[0,1]$.

  This phenomenon is a consequence of the fact that condition (3) of a series
  representation is shaped so as to constrain values only at the points where the
  sum of the dominating functions converges.
  If the base space has a point at which all elements of $L$ are forced to take
  the same value, then the Daniell extension becomes needlessly thin there.
\end{remark}

\begin{definition}
  Let $(X,L,I)$ and $(Y,M,J)$ be integration systems on sets $X$ and
  $Y$ respectively.
  A \textbf{morphism of integration systems}
  \[
    F:(X,L,I)\to(Y,M,J)
  \]
  is a morphism of Riesz systems $F:(X,L)\to(Y,M)$ such that
  \[
    J(F^*f)=I(f)
  \]
  holds for every $f\in L$.
\end{definition}

\begin{lemma}
  \begin{enumerate}
    \item For every integration system $(X,L,I)$ on a set $X$, the
      identity map $\operatorname{id}_X:X\to X$ induces a morphism of
      integration systems
      $\operatorname{id}_{(X,L,I)}:(X,L,I)\to (X,L,I)$.
    \item For any two morphisms of integration systems
      $F:(X,L,I)\to(Y,M,J)$ and $G:(Y,M,J)\to(Z,N,K)$, the map
      $\underline{F}\circ\underline{G}:Z\to X$ induces a morphism of integration
      systems $G\circ F:(X,L,I)\to(Z,N,K)$.
  \end{enumerate}
\end{lemma}
\begin{proof}
  We have already seen that the identity map induces a morphism of Riesz systems
  $\operatorname{id}_{(X,L)}$.
  Moreover $f\circ\underline{\operatorname{id}_X}=f\in L$ and
  $I(f\circ\underline{\operatorname{id}_X})=I(f)$ for every $f\in L$, so the
  identity map clearly induces a morphism of integration systems
  $\operatorname{id}_{(X,L,I)}:(X,L,I)\to (X,L,I)$.

  We have already seen that $\underline{F}\circ\underline{G}:Z\to X$ determines a
  morphism of Riesz systems from $(X,L)$ to $(Z,N)$.
  Moreover, by the hypothesis on morphisms of integration systems,
  \[
    K((G\circ F)^*(f))=K(G^*(F^*(f)))=J(F^*(f))=I(f) .
  \]
  Hence $\underline{F}\circ\underline{G}:Z\to X$ induces a morphism of
  integration systems $G\circ F:(X,L,I)\to(Z,N,K)$.
\end{proof}

\begin{definition}
  The category whose objects are integration systems and whose morphisms are
  morphisms of integration systems is written
  \[
    \mathbf{Int}
  \]
  and called the \textbf{category of integration systems}.
\end{definition}

\begin{remark}[on the direction of morphisms]\label{rem:opposite-convention}
  Note that the map of base spaces underlying a morphism
  $F:(X,L,I)\to(Y,M,J)$ is $\underline{F}:Y\to X$, so that the direction is
  reversed.
  In the language of measure theory, the condition
  $J(f\circ\underline{F})=I(f)$ says that $\underline{F}$ pushes $J$ forward to
  $I$, that is, that $\underline{F}$ is a measure-preserving map
  $(Y,J)\to(X,I)$.

  Hence \textbf{$\mathbf{Int}$ is the opposite of the usual category whose
  morphisms are measure-preserving maps}.
  This paper adopts this direction because pullback of functions is
  contravariant and we wish to treat the Daniell extension as a covariant
  functor.
  A reader who prefers to read everything in the usual convention should replace
  ``reflection'' below by ``coreflection'' and ``left adjoint'' by
  ``right adjoint''.
\end{remark}
